% !TEX root = ../main.tex
\chapter{The Geometry of Spacetime}\label{ch:geometry}

\begin{goals}
\begin{itemize}
\item Draw and read spacetime diagrams, worldlines and light cones.
\item Compute the interval between two events and say whether it is timelike,
  spacelike or null.
\item Compute the time a clock records along any path through spacetime.
\item Use a metric to measure intervals, and tell a coordinate effect from a
  real feature of the geometry.
\item Compute how fast clocks tick in a weak gravitational field and near a
  black hole.
\item Describe free fall as motion along a geodesic, and curvature as the
  source of tidal forces.
\end{itemize}
\end{goals}

A spacetime engineer designs geometries, so we need a precise language for
geometry. This chapter builds it in steps. We start with the flat spacetime of
special relativity and learn to measure it with the interval. We then write the
interval in a form, the \emph{metric}, that works in any coordinates and in
curved spacetime too. With the metric in hand, we can compute everything a
traveller would experience: the ticking of clocks, the paths of light, the
motion of falling objects, and the stretching forces called tides.

\section{Events, worldlines and units}\label{sec:events}

An \emph{event} is a place at an instant: a firecracker going off at a
particular point and time. We label events with four coordinates, a time $t$
and three positions $x$, $y$, $z$. A particle traces out a sequence of events
called its \emph{worldline}.

A \emph{spacetime diagram} plots worldlines with time running up the page and
one direction of space across it (\cref{fig:spacetime-diagram}). A particle at
rest has a vertical worldline. A moving particle has a tilted one, and the
faster it moves the more its worldline leans over.

It pays to measure time and distance in the same units. A light-second is the
distance light travels in one second, about \SI{3e8}{\metre}. If we measure
distance in light-seconds and time in seconds, or equivalently time in metres
(one metre of time is the time light takes to travel a metre, about
\SI{3.3}{\nano\second}), then the speed of light is exactly $c = 1$. Velocities
become pure numbers: a speed of $0.5$ means half the speed of light. From here
on we use these units, and restore factors of $c$ when we need a number in
everyday units. Light worldlines then run at $45^{\circ}$ on every spacetime
diagram.

\begin{figure}[htbp]
\centering
\begin{tikzpicture}[scale=1.0]
  \draw[-{Stealth}] (-3.2,0) -- (3.4,0) node[right] {$x$};
  \draw[-{Stealth}] (0,-0.4) -- (0,4.2) node[above] {$t$};
  \draw[very thick, worldline] (-2,0) -- (-2,3.9) node[above, font=\small, align=center] {at rest};
  \draw[very thick, examplegreen] (0,0) -- (1.6,3.9) node[above, font=\small] {moving, $v=0.4$};
  \draw[thick, lightcone!85!black, dashed] (0,0) -- (3.3,3.3) node[right, font=\small] {light};
  \draw[thick, lightcone!85!black, dashed] (0,0) -- (-3.2,3.2);
  \fill (0,0) circle (1.6pt) node[below right, font=\small] {event $A$};
\end{tikzpicture}
\caption{A spacetime diagram. Time runs up the page. A particle at rest has a
vertical worldline; a moving particle has a tilted one. With $c = 1$, light
travels along lines at $45^{\circ}$.}
\label{fig:spacetime-diagram}
\end{figure}

\section{The interval}\label{sec:interval}

In ordinary three-dimensional space, two observers who use differently
oriented axes disagree about the coordinates of two points but agree on the
distance between them. The squared distance
$\Delta x^{2} + \Delta y^{2} + \Delta z^{2}$ is the same in every set of
rotated axes.

Special relativity has an exact analogue. Observers moving relative to one
another disagree about time differences and about distances, yet all inertial
observers agree on the \emph{interval}
\begin{equation}
  \Delta s^{2} = -\Delta t^{2} + \Delta x^{2} + \Delta y^{2} + \Delta z^{2}.
  \label{eq:interval}
\end{equation}
The minus sign in front of the time difference is the whole difference between
the geometry of spacetime and the geometry of space. It divides pairs of events
into three kinds.

\begin{definition}[Kinds of interval]
The interval between two events is
\begin{itemize}
\item \emph{timelike} if $\Delta s^{2} < 0$: a clock can travel from one event
  to the other at less than the speed of light;
\item \emph{null} (or \emph{lightlike}) if $\Delta s^{2} = 0$: a light signal
  can travel from one to the other;
\item \emph{spacelike} if $\Delta s^{2} > 0$: nothing, not even light, can
  travel from one to the other.
\end{itemize}
\end{definition}

The events that light can reach from an event $A$ form a cone, the
\emph{future light cone} of $A$ (\cref{fig:light-cone}). Events inside it can
be influenced by what happens at $A$; events outside it cannot. The events
from which light can reach $A$ form its past light cone. The light cones of all
events together are the \emph{causal structure} of spacetime: they tell you
what can affect what. Every question about faster-than-light travel is, at
bottom, a question about light cones.

\begin{figure}[htbp]
\centering
\begin{tikzpicture}[scale=0.95]
  \fill[lightcone!18] (0,0) -- (-2.2,2.2) arc (180:360:2.2 and 0.45) -- cycle;
  \draw[lightcone!80!black, thick] (0,0) -- (-2.2,2.2);
  \draw[lightcone!80!black, thick] (0,0) -- (2.2,2.2);
  \draw[lightcone!80!black, thick] (0,2.2) ellipse (2.2 and 0.45);
  \fill[lightcone!10] (0,0) -- (-2.2,-2.2) arc (180:360:2.2 and 0.45) -- cycle;
  \draw[lightcone!80!black, thick] (0,0) -- (-2.2,-2.2);
  \draw[lightcone!80!black, thick] (0,0) -- (2.2,-2.2);
  \draw[lightcone!80!black, thick, dashed] (0,-2.2) ellipse (2.2 and 0.45);
  \draw[-{Stealth}] (0,-2.9) -- (0,3.1) node[above] {$t$};
  \fill (0,0) circle (1.7pt) node[right, font=\small] {$\;A$};
  \node[font=\small, align=center] at (0,1.5) {future};
  \node[font=\small, align=center] at (0,-1.5) {past};
  \node[font=\small, align=center] at (3.6,0.2) {elsewhere\\(spacelike from $A$)};
  \fill[examplegreen] (0.5,1.3) circle (1.5pt) node[right, font=\footnotesize] {can be influenced by $A$};
  \fill[deeperviolet] (2.6,-0.9) circle (1.5pt) node[right, font=\footnotesize] {cannot};
\end{tikzpicture}
\caption{The light cone of an event $A$, drawn with one dimension of space
suppressed. Events inside the future cone can be influenced by $A$. Events
outside both cones are spacelike from $A$, and no signal connects them to it.}
\label{fig:light-cone}
\end{figure}

\begin{checkpoint}
Event $B$ happens \SI{3}{\second} after event $A$ and 4~light-seconds
away from it. Event $C$ happens \SI{5}{\second} after $A$ and
4~light-seconds away. Classify the intervals $AB$ and $AC$. Can
anything that happens at $A$ affect $B$? Can it affect $C$?
\end{checkpoint}

\section{Proper time}\label{sec:proper-time}

A clock carried along a worldline measures the \emph{proper time} $\tau$ along
it. For a clock at rest, proper time is just coordinate time. For a moving
clock, the interval does the bookkeeping. Along a short piece of a timelike
worldline the interval is negative, and the clock records
\begin{equation}
  \dd\tau^{2} = -\dd s^{2} = \dd t^{2} - \dd x^{2} - \dd y^{2} - \dd z^{2}.
  \label{eq:proper-time}
\end{equation}
Dividing by $\dd t^{2}$ gives $\dd\tau/\dd t = \sqrt{1 - v^{2}}$, where $v$ is the
clock's speed. The proper time along a whole worldline is the sum of these
pieces:
\begin{equation}
  \tau = \int \sqrt{1 - v(t)^{2}}\;\dd t .
  \label{eq:proper-time-integral}
\end{equation}
This is the time dilation of special relativity written as a geometric length:
proper time is to a worldline what length is to a curve in space.

\begin{example}{A trip to the nearest star}\label{ex:twin}
Proxima Centauri is about 4.24 light-years away. A traveller leaves Earth at
$v = 0.8$, reaches the star, turns around at once, and returns at the same
speed. How much time passes on Earth, and on the traveller's clock?

\solution In the Earth's frame the trip covers $2 \times 4.24 = 8.48$
light-years at speed $0.8$, so it takes
\begin{equation*}
  \Delta t = \frac{8.48}{0.8}\ \text{years} = 10.6\ \text{years}.
\end{equation*}
The traveller moves at constant speed on each leg, so
\cref{eq:proper-time-integral} gives
\begin{equation*}
  \Delta\tau = \sqrt{1 - 0.8^{2}}\;\Delta t = 0.6 \times 10.6\ \text{years}
  = 6.36\ \text{years}.
\end{equation*}
\discussion The traveller returns 4.2 years younger than a twin who stayed
home. The trip is not symmetric: the Earth twin's worldline is straight, while
the traveller's worldline has a corner at the turnaround. In the geometry of
spacetime, the straight worldline between two events is the one with the
\emph{longest} proper time, which is the reverse of ordinary geometry, where the
straight line is the shortest. The minus sign in \cref{eq:interval} is
responsible.
\end{example}

\begin{keyidea}
Between two events, the straight worldline of a freely moving clock records
the longest proper time. Every detour, every acceleration, shortens the time a
clock records.
\end{keyidea}

\section{The metric}\label{sec:metric}

\Cref{eq:interval} holds for inertial observers using ordinary Cartesian
coordinates. We need a version that works in any coordinates, and later in
curved spacetime. The idea is to write down the interval between
\emph{neighbouring} events, as a formula in whatever coordinates we are using.

\subsection{Index notation}

We label the four coordinates with an index that runs from 0 to 3:
$x^{0} = t$, $x^{1} = x$, $x^{2} = y$, $x^{3} = z$. The raised index is a label,
not a power. Greek letters $\mu, \nu, \dots$ stand for any of the four values.
The interval between neighbouring events can then be written
\begin{equation}
  \dd s^{2} = \sum_{\mu=0}^{3}\sum_{\nu=0}^{3} g_{\mu\nu}\,\dd x^{\mu}\,\dd x^{\nu}
  \equiv g_{\mu\nu}\,\dd x^{\mu}\,\dd x^{\nu}.
  \label{eq:metric}
\end{equation}
The sixteen numbers $g_{\mu\nu}$ form a symmetric $4\times4$ matrix called the
\emph{metric}. In the second form we have used the \emph{summation convention}:
an index that appears once up and once down in a term is summed from 0 to 3.
We use the convention from now on.

For flat spacetime in Cartesian coordinates, comparing with
\cref{eq:interval} gives the Minkowski metric
\begin{equation}
  g_{\mu\nu} = \eta_{\mu\nu} =
  \begin{pmatrix} -1 & 0 & 0 & 0\\ 0 & 1 & 0 & 0\\ 0 & 0 & 1 & 0\\ 0 & 0 & 0 & 1\end{pmatrix}.
\end{equation}

\subsection{Coordinates are labels}

The same flat spacetime looks different in other coordinates. In spherical
coordinates the interval reads
\begin{equation}
  \dd s^{2} = -\dd t^{2} + \dd r^{2} + r^{2}\bigl(\dd\theta^{2} + \sin^{2}\theta\,\dd\phi^{2}\bigr).
  \label{eq:flat-spherical}
\end{equation}
The metric components now depend on position, yet nothing physical has
changed. Coordinates are labels we attach to events, and a good deal of the
skill of working with metrics lies in telling which features of a metric belong
to the labels and which belong to the geometry. The next example shows how
dramatic a pure coordinate effect can look.

\begin{example}{Flowing coordinates}\label{ex:flow}
Consider the metric
\begin{equation}
  \dd s^{2} = -\dd t^{2} + \bigl(\dd x - v_{0}\,\dd t\bigr)^{2} + \dd y^{2} + \dd z^{2},
  \label{eq:flow-metric}
\end{equation}
with $v_{0}$ a constant. Find the coordinate speed $\dd x/\dd t$ of light moving
along $x$. Then show that this metric describes flat spacetime.

\solution Light has $\dd s^{2} = 0$. Along $x$ we set $\dd y = \dd z = 0$ and find
$(\dd x - v_{0}\dd t)^{2} = \dd t^{2}$, so $\dd x - v_{0}\,\dd t = \pm\dd t$ and
\begin{equation*}
  \frac{\dd x}{\dd t} = v_{0} \pm 1.
\end{equation*}
For $v_{0} = 0.5$, light moves at $+1.5$ in one direction and $-0.5$ in the
other, as measured in these coordinates.

Now change coordinates to $x' = x - v_{0}t$, keeping $t$, $y$ and $z$. Then
$\dd x' = \dd x - v_{0}\,\dd t$, and the metric becomes
$\dd s^{2} = -\dd t^{2} + \dd x'^{2} + \dd y^{2} + \dd z^{2}$, the Minkowski metric.

\discussion This is flat spacetime described by an observer who labels points
with a grid that slides along $x$ at speed $v_{0}$ relative to inertial
observers. In the sliding grid, light seems to move faster than 1 in one
direction. No object ever outruns a light signal sent beside it: the speeds
$v_{0} \pm 1$ describe how the labels move, not how anything moves relative to
its surroundings. Keep this example in mind. The metrics of warp drives have
exactly this form, with $v_{0}$ replaced by a function that varies from place
to place, and in that case the effect is no longer only a matter of labels.
\end{example}

The form \cref{eq:flow-metric} gives us a picture that will run through this
book: space \emph{flowing} past a grid of coordinates. Observers who are carried
along by the flow feel no force. The flow speed can take any value, and where
it varies from place to place the geometry is genuinely curved.

\begin{checkpoint}
In \cref{ex:flow}, suppose $v_{0} = 2$. In which direction can light move, as
measured in the sliding coordinates? Does this mean something is moving faster
than light?
\end{checkpoint}

\subsection{Four-velocity}

The tangent to a worldline, measured per unit proper time, is the
\emph{four-velocity} $u^{\mu} = \dd x^{\mu}/\dd\tau$. Dividing
\cref{eq:metric} by $\dd\tau^{2} = -\dd s^{2}$ gives the normalization
\begin{equation}
  g_{\mu\nu}\,u^{\mu}u^{\nu} = -1.
  \label{eq:four-velocity}
\end{equation}
For a clock at rest in flat spacetime, $u^{\mu} = (1,0,0,0)$. For a clock moving
at speed $v$ along $x$, $u^{\mu} = \gamma\,(1, v, 0, 0)$ with
$\gamma = 1/\sqrt{1-v^{2}}$. We will use four-velocities in
\cref{ch:matter-geometry} to say what an observer measures.

\section{Clocks in curved spacetime}\label{sec:clocks}

In curved spacetime no choice of coordinates puts the metric into the
Minkowski form everywhere at once. The metric still gives the interval between
neighbouring events through \cref{eq:metric}, and the proper time along a
worldline is still $\tau = \int\sqrt{-\dd s^{2}}$. Everything we did above
carries over; only the metric changes.

\subsection{Weak gravity}

Around a planet or a star the gravitational field is weak, and the metric is
very close to Minkowski. To an excellent approximation,
\begin{equation}
  \dd s^{2} = -(1 + 2\Phi)\,\dd t^{2} + (1 - 2\Phi)\bigl(\dd x^{2} + \dd y^{2} + \dd z^{2}\bigr),
  \label{eq:weak-field}
\end{equation}
where $\Phi$ is the Newtonian gravitational potential divided by $c^{2}$. For a
spherical mass $M$, $\Phi = -GM/(c^{2}r)$, which is small and negative:
$\Phi \approx -7\times10^{-10}$ at the Earth's surface. The approximation holds
whenever $|\Phi| \ll 1$.

A clock moving slowly at speed $v$ through this spacetime records, from
\cref{eq:weak-field},
\begin{equation}
  \frac{\dd\tau}{\dd t} = \sqrt{1 + 2\Phi - (1 - 2\Phi)\,v^{2}}
  \approx 1 + \Phi - \tfrac{1}{2}v^{2},
  \label{eq:clock-rate}
\end{equation}
keeping only the leading terms in the small quantities $\Phi$ and $v^{2}$. The
two corrections have a clear meaning. A clock deeper in the potential (more
negative $\Phi$) runs slower, and a moving clock runs slower.

\begin{example}{Clocks on the ground and in orbit}\label{ex:gps}
Navigation satellites orbit at a radius $r_{\mathrm{s}} = \SI{26560}{\kilo\metre}$
from the Earth's centre. Take $GM = \SI{3.986e14}{\metre\cubed\per\second\squared}$
for the Earth and $R = \SI{6371}{\kilo\metre}$ for its radius, and ignore the
Earth's rotation. By how much does a satellite clock gain or lose each day,
compared with a clock on the ground?

\solution The ground clock is at rest at $r = R$, so its rate is
$1 + \Phi(R)$. The satellite clock moves on a circular orbit at speed
$v = \sqrt{GM/r_{\mathrm{s}}}$, about \SI{3.87}{\kilo\metre\per\second}. The
difference in rates is
\begin{equation*}
  \frac{\dd\tau_{\mathrm{s}}}{\dd t} - \frac{\dd\tau_{\mathrm{g}}}{\dd t}
  = \underbrace{\frac{GM}{c^{2}}\left(\frac{1}{R} - \frac{1}{r_{\mathrm{s}}}\right)}_{\text{altitude}}
  \; - \; \underbrace{\frac{v^{2}}{2c^{2}}}_{\text{speed}} .
\end{equation*}
Numerically the altitude term is $5.29\times10^{-10}$ and the speed term is
$8.35\times10^{-11}$. Over one day, \SI{86400}{\second}, the satellite clock
gains \SI{45.7}{\micro\second} from its altitude and loses
\SI{7.2}{\micro\second} from its speed, a net gain of about
\SI{38.5}{\micro\second} per day.

\discussion Light travels \SI{11.5}{\kilo\metre} in \SI{38.5}{\micro\second},
so an uncorrected clock would spoil positions by kilometres within a day. The
full treatment, including the Earth's rotation and shape, is given by
\textcite{Ashby2003}. This is the calculation behind the numbers quoted in
\cref{sec:gps}.
\end{example}

\subsection{Near a black hole}

For a non-rotating mass $M$ with no weak-field approximation at all, the
geometry outside the mass is described by the \emph{Schwarzschild metric}:
\begin{equation}
  \dd s^{2} = -\left(1 - \frac{r_{\mathrm{s}}}{r}\right)\dd t^{2}
  + \frac{\dd r^{2}}{1 - r_{\mathrm{s}}/r}
  + r^{2}\bigl(\dd\theta^{2} + \sin^{2}\theta\,\dd\phi^{2}\bigr),
  \qquad r_{\mathrm{s}} = \frac{2GM}{c^{2}} .
  \label{eq:schwarzschild}
\end{equation}
Here $r$ is defined so that a sphere at radius $r$ has area $4\pi r^{2}$. A
clock held at fixed $r$ records $\dd\tau/\dd t = \sqrt{1 - r_{\mathrm{s}}/r}$
relative to a clock far away. For the Sun, $r_{\mathrm{s}} \approx
\SI{2.95}{\kilo\metre}$, so the effect at its surface is tiny. For a neutron star
of 1.4 solar masses and radius \SI{12}{\kilo\metre}, the surface clock runs at
0.81 of the rate of a distant clock (\cref{prob:neutron-star}).

At $r = r_{\mathrm{s}}$ the clock rate would fall to zero. No clock can be held
at rest there: $r = r_{\mathrm{s}}$ is the \emph{event horizon} of a black hole.
The metric \cref{eq:schwarzschild} also misbehaves at the horizon, where
$g_{rr}$ becomes infinite. That part is a coordinate effect, as another way of
writing the same geometry shows.

\subsection{Space as a flowing river}

Introduce a new time coordinate $T$, the time on the clocks of observers who
fall inward from rest far away. In these coordinates the Schwarzschild geometry
takes the Painlevé--Gullstrand form (in units with $c = 1$):
\begin{equation}
  \dd s^{2} = -\dd T^{2} + \left(\dd r + \sqrt{\frac{r_{\mathrm{s}}}{r}}\,\dd T\right)^{2}
  + r^{2}\bigl(\dd\theta^{2} + \sin^{2}\theta\,\dd\phi^{2}\bigr).
  \label{eq:river}
\end{equation}
Compare this with \cref{eq:flow-metric}. It has the same form, with space
flowing \emph{inward} at the speed $\sqrt{r_{\mathrm{s}}/r}$, which grows as
$r$ shrinks. Nothing is singular at $r = r_{\mathrm{s}}$ any more.

\begin{example}{Where light stands still}\label{ex:horizon}
Use \cref{eq:river} to find the coordinate speed $\dd r/\dd T$ of light moving
radially, and locate the place where outgoing light makes no progress.

\solution For radial light, $\dd s^{2} = 0$ with $\dd\theta = \dd\phi = 0$, so
$\dd r + \sqrt{r_{\mathrm{s}}/r}\,\dd T = \pm\dd T$, or
\begin{equation*}
  \frac{\dd r}{\dd T} = -\sqrt{\frac{r_{\mathrm{s}}}{r}} \pm 1 .
\end{equation*}
Ingoing light ($-$ sign) always moves inward. Outgoing light ($+$ sign) moves
outward only where $\sqrt{r_{\mathrm{s}}/r} < 1$, that is, where $r > r_{\mathrm{s}}$.
At $r = r_{\mathrm{s}}$ it stands still, and inside it is swept inward.

\discussion This is the horizon, seen as a place where space flows inward faster
than light can move through it. Light always moves at speed 1 relative to the
space around it; the horizon forms where the flow itself exceeds 1. A warp drive
uses the same mechanism in reverse, carrying a region along by a flow of space.
It also inherits the same danger: wherever the flow outruns light, horizons can
form. You will meet those horizons again in Part~III.
\end{example}

\begin{figure}[htbp]
\centering
\begin{tikzpicture}[scale=0.95]
  \draw[-{Stealth}] (0,0) -- (10.2,0) node[right] {$r$};
  \draw[thick, deeperviolet, dashed] (3,-1.5) -- (3,1.75);
  \node[font=\small, text=deeperviolet] at (3,2.05) {horizon, $r = r_{\mathrm s}$};
  \draw[-{Stealth[length=2mm]}, thick, keygold!80!black] (1.912,1.0) -- (0.488,1.0);
  \draw[-{Stealth[length=2mm]}, thick, keygold!80!black] (2.725,1.0) -- (1.675,1.0);
  \draw[-{Stealth[length=2mm]}, thick, keygold!80!black] (3.450,1.0) -- (2.550,1.0);
  \draw[-{Stealth[length=2mm]}, thick, keygold!80!black] (4.867,1.0) -- (4.133,1.0);
  \draw[-{Stealth[length=2mm]}, thick, keygold!80!black] (6.806,1.0) -- (6.194,1.0);
  \draw[-{Stealth[length=2mm]}, thick, keygold!80!black] (9.260,1.0) -- (8.740,1.0);
  \draw[-{Stealth[length=2mm]}, thick, examplegreen] (1.462,-0.9) -- (0.938,-0.9);
  \draw[-{Stealth[length=2mm]}, thick, examplegreen] (2.275,-0.9) -- (2.125,-0.9);
  \fill[examplegreen] (3.000,-0.9) circle (2.2pt);
  \draw[-{Stealth[length=2mm]}, thick, examplegreen] (4.417,-0.9) -- (4.583,-0.9);
  \draw[-{Stealth[length=2mm]}, thick, examplegreen] (6.356,-0.9) -- (6.644,-0.9);
  \draw[-{Stealth[length=2mm]}, thick, examplegreen] (8.810,-0.9) -- (9.190,-0.9);
  \node[font=\small, text=keygold!70!black, anchor=west] at (10.4,1.0) {flow of space};
  \node[font=\small, text=examplegreen, anchor=west, align=left] at (10.4,-0.9) {outgoing light,\\net motion};
\end{tikzpicture}
\caption{The Schwarzschild geometry pictured as space flowing inward
(\cref{eq:river}). Top: the inward flow speed $\sqrt{r_{\mathrm s}/r}$ at six radii, drawn
to scale. Bottom: the net coordinate velocity $1 - \sqrt{r_{\mathrm s}/r}$ of outgoing
light at the same radii. Outside the horizon light makes headway; at the horizon
it stands still; inside, the flow sweeps it inward.}
\label{fig:river}
\end{figure}

\section{Free fall and geodesics}\label{sec:geodesics}

In \cref{ex:twin} we found that the straight worldline between two events has
the longest proper time. General relativity turns this into the law of motion
for anything on which no non-gravitational force acts:

\begin{keyidea}[free fall]
A freely falling body follows a worldline of extremal proper time, a
\emph{geodesic} of spacetime. Gravity is not a force that deflects it; it is
the curvature that bends what \enquote{straight} means.
\end{keyidea}

Making the proper time extremal leads, after a calculation you can do in
\cref{prob:geodesic}, to the \emph{geodesic equation}
\begin{equation}
  \frac{\dd^{2}x^{\mu}}{\dd\tau^{2}}
  + \Gamma^{\mu}{}_{\alpha\beta}\,\frac{\dd x^{\alpha}}{\dd\tau}\frac{\dd x^{\beta}}{\dd\tau} = 0,
  \label{eq:geodesic}
\end{equation}
where the \emph{Christoffel symbols} are built from first derivatives of the
metric:
\begin{equation}
  \Gamma^{\mu}{}_{\alpha\beta} = \tfrac{1}{2}\,g^{\mu\nu}
  \left(\partial_{\alpha}g_{\nu\beta} + \partial_{\beta}g_{\nu\alpha} - \partial_{\nu}g_{\alpha\beta}\right).
  \label{eq:christoffel}
\end{equation}
Here $\partial_{\alpha}$ means $\partial/\partial x^{\alpha}$, and $g^{\mu\nu}$ is
the inverse of the matrix $g_{\mu\nu}$. In flat spacetime with Cartesian
coordinates every $\Gamma$ vanishes, and \cref{eq:geodesic} says that free
bodies move in straight lines at constant speed.

\begin{example}{Newton's law of gravity from geometry}\label{ex:newton}
Show that for a slowly moving body in the weak-field metric
\cref{eq:weak-field}, the geodesic equation reduces to Newton's law
$\dd^{2}x^{i}/\dd t^{2} = -\partial_{i}\Phi$, for a potential that does not
change in time.

\solution A slow body has $\dd x^{i}/\dd\tau \ll \dd t/\dd\tau$, so in
\cref{eq:geodesic} the term with $\alpha = \beta = 0$ dominates:
\begin{equation*}
  \frac{\dd^{2}x^{i}}{\dd\tau^{2}} \approx -\Gamma^{i}{}_{00}\left(\frac{\dd t}{\dd\tau}\right)^{2}.
\end{equation*}
With $g_{00} = -(1+2\Phi)$ independent of time, \cref{eq:christoffel} gives
$\Gamma^{i}{}_{00} = -\tfrac{1}{2}g^{ii}\partial_{i}g_{00} \approx \partial_{i}\Phi$ to
first order. Since $\dd t/\dd\tau \approx 1$ for a slow body,
\begin{equation*}
  \frac{\dd^{2}x^{i}}{\dd t^{2}} \approx -\partial_{i}\Phi .
\end{equation*}
\discussion Newton's law of gravity emerges from the $tt$ part of the metric
alone, that is, from the way the rate of clocks changes from place to place.
Objects fall toward regions where clocks run slower. This is one of the most
important intuitions in the book: the rate of clocks, and how it varies across
space, is itself a gravitational field.
\end{example}

\section{Curvature and tides}\label{sec:tides}

A single freely falling observer feels nothing: in a small enough laboratory,
free fall is indistinguishable from floating in empty space. This is the
\emph{equivalence principle}, and it is why the metric can always be made to
look like Minkowski at any one event by a suitable choice of coordinates.
Gravity reveals itself when you compare \emph{nearby} free-fall paths.

Drop two balls side by side above the Earth. Both fall toward the Earth's
centre, so their paths converge and the balls drift together. Drop two balls one
above the other, and the lower one, closer to the Earth, accelerates slightly
more; they drift apart. These relative accelerations are \emph{tides}, and no
choice of coordinates can remove them. They are the true signature of curvature.

\subsection{The Riemann tensor}

In general relativity the relative acceleration of two nearby free-fall paths
separated by a small displacement $\xi^{\mu}$ is
\begin{equation}
  \frac{\dd^{2}\xi^{\mu}}{\dd\tau^{2}} = -R^{\mu}{}_{\alpha\nu\beta}\,u^{\alpha}\xi^{\nu}u^{\beta},
  \label{eq:geodesic-deviation}
\end{equation}
where $u^{\mu}$ is the four-velocity of the falling bodies and
$R^{\mu}{}_{\alpha\nu\beta}$ is the \emph{Riemann curvature tensor}:
\begin{equation}
  R^{\mu}{}_{\alpha\nu\beta} = \partial_{\nu}\Gamma^{\mu}{}_{\alpha\beta}
  - \partial_{\beta}\Gamma^{\mu}{}_{\alpha\nu}
  + \Gamma^{\mu}{}_{\nu\lambda}\Gamma^{\lambda}{}_{\alpha\beta}
  - \Gamma^{\mu}{}_{\beta\lambda}\Gamma^{\lambda}{}_{\alpha\nu}.
  \label{eq:riemann}
\end{equation}
The Riemann tensor vanishes everywhere if and only if spacetime is flat. When
\cref{ex:flow} asked whether a metric describes flat spacetime, computing its
Riemann tensor would have answered the question directly, without having to
guess the right change of coordinates.

In the weak-field limit, the tidal part of the Riemann tensor becomes the
familiar Newtonian tidal tensor, $R^{i}{}_{0j0} \approx \partial_{i}\partial_{j}\Phi$.
Near the Earth the tidal acceleration across a person \SI{2}{\metre} tall is
about $2GM L/r^{3} \approx \SI{6e-6}{\metre\per\second\squared}$, far too small to
notice. Near a black hole of a few solar masses it would tear a person apart.
Tides are one of the first things a spacetime engineer checks on behalf of a
traveller.

\subsection{Ricci and Einstein}

Two combinations of the Riemann tensor carry the information Einstein's
equation needs. The \emph{Ricci tensor} and the \emph{Ricci scalar} are
\begin{equation}
  R_{\alpha\beta} = R^{\mu}{}_{\alpha\mu\beta}, \qquad R = g^{\alpha\beta}R_{\alpha\beta},
\end{equation}
and the \emph{Einstein tensor} is
\begin{equation}
  G_{\alpha\beta} = R_{\alpha\beta} - \tfrac{1}{2}R\,g_{\alpha\beta}.
  \label{eq:einstein-tensor}
\end{equation}
The Einstein tensor is the \enquote{curvature of spacetime} on the left side of
Einstein's equation. \Cref{ch:matter-geometry} explains why this particular
combination appears, and uses it to compute what a geometry demands.

\begin{deeper}{tensors and coordinate independence}
The objects in this chapter are \emph{tensors}: quantities whose components
change in a definite way under a change of coordinates, so that equations
between tensors hold in every coordinate system if they hold in one. The metric
is a symmetric tensor of type $(0,2)$; the Riemann tensor has type $(1,3)$.
The Christoffel symbols are \emph{not} tensor components, which is why they can
vanish at a point in one set of coordinates and not in another. Spacetime is
modelled as a four-dimensional Lorentzian manifold, and all the curvature
quantities here follow the sign conventions of \textcite{Misner1973}. Wald's
text develops the same material in coordinate-free, abstract-index form
\parencite{Wald1984}.
\end{deeper}

\begin{chaptersummary}
\item With $c = 1$, time and distance share units and light moves at
  $45^{\circ}$ on a spacetime diagram.
\item The interval $\Delta s^{2} = -\Delta t^{2} + \Delta x^{2} + \Delta y^{2} + \Delta z^{2}$
  is the same for all inertial observers. It is timelike, null or spacelike,
  and light cones divide spacetime into what an event can and cannot
  influence.
\item A clock measures proper time $\tau = \int\sqrt{-\dd s^{2}}$ along its
  worldline. Between two events, the free-fall worldline records the longest
  proper time.
\item The metric $g_{\mu\nu}$ gives the interval between neighbouring events in
  any coordinates. Coordinate effects can be dramatic, as in flowing
  coordinates, where light's coordinate speed differs from 1.
\item In a weak field a clock runs at $\dd\tau/\dd t \approx 1 + \Phi - v^{2}/2$. The
  GPS satellites gain about \SI{38}{\micro\second} per day for this reason.
\item The Schwarzschild geometry can be pictured as space flowing inward; its
  horizon is where the flow reaches the speed of light.
\item Free bodies follow geodesics, and the rate of clocks from place to place
  acts as a gravitational field. The Riemann tensor measures tides, the true
  signature of curvature; the Ricci and Einstein tensors are built from it.
\end{chaptersummary}

\begin{checkanswers}
\item[2.1] $AB$: $\Delta s^{2} = -9 + 16 = +7$, spacelike, so nothing at $A$ can
  affect $B$. $AC$: $\Delta s^{2} = -25 + 16 = -9$, timelike, so a signal slower
  than light can carry an influence from $A$ to $C$.
\item[2.2] With $v_{0} = 2$, both solutions $\dd x/\dd t = 2 \pm 1$ are positive:
  light moves only toward $+x$ in the sliding coordinates. Nothing outruns
  light. The grid slides at twice the speed of light relative to inertial
  observers, which is allowed because a grid of labels is not an object.
\end{checkanswers}

\begin{problems}
\item \diff{1} Two events are separated by $\Delta t = \SI{2}{\second}$ and
  $\Delta x = 1.5$~light-seconds. Find the interval and classify it. What
  proper time does a clock record if it travels between them at constant
  velocity?
\item \diff{1} A clock travels at $v = 0.99$ for one year of Earth time. How much
  time does it record?
\item \diff{2} \textbf{Accelerating twin.} Repeat \cref{ex:twin} with a more
  realistic trip: the traveller accelerates uniformly from rest to $0.8$ over
  the first year of Earth time, cruises, and decelerates symmetrically at the
  star and on the way home. Set up the proper-time integral and evaluate it
  numerically.
\item \diff{1} \textbf{Everest.} By what fraction does a clock on the summit of
  Mount Everest (\SI{8.85}{\kilo\metre}) run faster than one at sea level?
  How many nanoseconds does it gain per year?
\item \diff{2}\label{prob:neutron-star} \textbf{Neutron star.} Show that a clock on
  the surface of a neutron star of mass $1.4\,\Msun$ and radius
  \SI{12}{\kilo\metre} runs at 0.81 of the rate of a distant clock. Use
  $GM_{\odot}/c^{2} = \SI{1.477}{\kilo\metre}$.
\item \diff{2} \textbf{Polar coordinates.} For the flat two-dimensional metric
  $\dd s^{2} = \dd r^{2} + r^{2}\dd\phi^{2}$, compute all the Christoffel symbols
  from \cref{eq:christoffel}. Show that the geodesic equation has straight
  lines as solutions.
\item \diff{2} \textbf{Flowing coordinates, again.} Compute the Christoffel
  symbols of the metric \cref{eq:flow-metric} and show that the Riemann tensor
  vanishes. Then repeat with $v_{0}$ replaced by a function $v(y)$, and show
  that the Riemann tensor no longer vanishes.
\item \diff{2}\label{prob:geodesic} \textbf{Deriving the geodesic equation.}
  Write the proper time along a worldline as $\tau = \int
  \sqrt{-g_{\mu\nu}\dot x^{\mu}\dot x^{\nu}}\,\dd\lambda$, where the dot is
  $\dd/\dd\lambda$. Apply the Euler--Lagrange equations and choose
  $\lambda = \tau$ to obtain \cref{eq:geodesic}.
\item \diff{2} \textbf{The river.} Show that \cref{eq:river} turns into
  \cref{eq:schwarzschild} under the change of time coordinate
  $\dd T = \dd t + \frac{\sqrt{r_{\mathrm{s}}/r}}{1 - r_{\mathrm{s}}/r}\,\dd r$.
\item \diff{3} \textbf{Tides near a black hole.} Estimate the stretching
  acceleration across a \SI{2}{\metre} tall person falling feet first into a
  black hole of $10\,\Msun$ at the moment they cross the horizon, and again for a
  black hole of $10^{9}\,\Msun$. Which one could a traveller cross unharmed, and
  why does the answer depend on the mass the way it does?
\end{problems}

\begin{furtherreading}
\item[\textcite{Hartle2003}] A physics-first introduction to general relativity
  that builds intuition for metrics, clocks and geodesics before the formalism.
\item[\textcite{Schutz2009}] A careful, gentle development of spacetime
  diagrams, tensors and curvature, well suited to self-study.
\item[\textcite{Carroll2004}] A graduate-level introduction that covers the
  geometry of this chapter with full mathematical care.
\item[\textcite{Misner1973}] The encyclopedic classic, with deep geometric
  intuition; its conventions are the ones used in this book.
\end{furtherreading}
